\section{讨论}
\label{sec:discussion}

本文提出了 DreamerV2，这是一种基于模型的智能体。在 Atari 200M 基准上，它仅通过单独训练的世界模型在潜空间中的预测来学习行为，性能达到人类水平。DreamerV2 在相同计算预算和训练时间下超过了顶级的单 GPU 无模型智能体 Rainbow。我们以 Dreamer 智能体\citep{hafner2019dreamer}为基础，并通过实验确认，学习分类潜空间以及使用 KL 平衡能够提升智能体性能。DreamerV2 依靠图像信息学习通用表征；即使这些表征并不是专门为预测回报而学习的，性能也不会受到明显影响。

DreamerV2 表明，在最具竞争性的强化学习基准之一上，基于模型的 RL 可以超过顶级无模型算法，尽管后者已经积累了多年的研究和工程优化。除了在单个任务上取得较强结果之外，世界模型还可用于高效迁移、多任务学习、物理机器人上的样本高效学习，以及基于不确定性估计的全局探索。

\paragraph{鸣谢}
我们感谢匿名审稿人的反馈，也感谢 Nick Rhinehart 就分类潜变量潜在优势所进行的深入讨论。
